% PAGES: 211-212

\begin{lemma}
Let $\Sigma_{\bfx}$ and $\Omega_{\bfx}$ be the covariance and the correlation matrix
of $n$ nonconstant random variables $X_1, X_2, \ldots, X_n$. Let $\sigma_i \neq 0$ be
the standard deviation of $X_i$, $i = 1:n$.

(i) The covariance matrix $\Sigma_{\bfx}$ and the correlation matrix $\Omega_{\bfx}$
are related as follows:
\begin{equation}
\Sigma_{\bfx} \;=\; D_{\sigma_{\bfx}}\Omega_{\bfx}D_{\sigma_{\bfx}},
\end{equation}
where $D_{\sigma_{\bfx}} = \diag(\sigma_i)_{i=1:n}$.

(ii) The correlation matrix $\Omega_{\bfx}$ is uniquely determined in terms of the
covariance matrix $\Sigma_{\bfx}$ as follows:
\begin{equation}
\Omega_{\bfx} \;=\; (D_{\sigma_{\bfx}})^{-1}\Sigma_{\bfx}(D_{\sigma_{\bfx}})^{-1},
\end{equation}
where $(D_{\sigma_{\bfx}})^{-1} = \diag\!\left(\frac{1}{\sigma_i}\right)_{i=1:n}$.
\end{lemma}

\begin{proof}
(i) If $D_1 = \diag\!\left(d_j^{(1)}\right)_{j=1:n}$ and $D_2 =
\diag\!\left(d_k^{(2)}\right)_{k=1:n}$ are diagonal matrices of size $n$, it follows
from (1.90) of Lemma 1.10 that
\begin{equation}
(D_1AD_2)(j,k) \;=\; d_j^{(1)}d_k^{(2)}A(j,k), \quad \forall\, 1\le j,k\le n.
\end{equation}
Using (7.18) for $A = \Omega_{\bfx}$ and for $D_1 = D_2 = D_{\sigma_{\bfx}} =
\diag(\sigma_i)_{i=1:n}$, we find that
\begin{align}
(D_{\sigma_{\bfx}}\Omega_{\bfx}D_{\sigma_{\bfx}})(j,k) &= \sigma_j\sigma_k\Omega_{\bfx}(j,k) \notag \\
&= \sigma_j\sigma_k\corr(X_j,X_k) \;=\; \cov(X_j,X_k) \notag \\
&= \Sigma_{\bfx}(j,k), \quad \forall\, 1\le j,k\le n,
\end{align}
since $\cov(X_j,X_k) = \sigma_j\sigma_k\corr(X_j,X_k)$; cf. (7.2). From (7.19), we
conclude that $\Sigma_{\bfx} = D_{\sigma_{\bfx}}\Omega_{\bfx}D_{\sigma_{\bfx}}$.

(ii) Since $\sigma_i \neq 0$ for all $i = 1:n$, the matrix $D_{\sigma_{\bfx}} =
\diag(\sigma_i)_{i=1:n}$ is nonsingular and its inverse is $(D_{\sigma_{\bfx}})^{-1} =
\diag\!\left(\frac{1}{\sigma_i}\right)_{i=1:n}$; see (1.94). Then, by multiplying
(7.16) from the left and from the right by $(D_{\sigma_{\bfx}})^{-1}$, we obtain that
\[
(D_{\sigma_{\bfx}})^{-1}\Sigma_{\bfx}(D_{\sigma_{\bfx}})^{-1}
\;=\; \big((D_{\sigma_{\bfx}})^{-1}D_{\sigma_{\bfx}}\big)\Omega_{\bfx}\big(D_{\sigma_{\bfx}}(D_{\sigma_{\bfx}})^{-1}\big)
\;=\; \Omega_{\bfx}.
\]
\end{proof}

Let $\Sigma_{\bfx}$ and $\Omega_{\bfx}$ be the covariance matrix and the correlation
matrix of two random variables $X_1$ and $X_2$ given by (7.12) and (7.13),
respectively. Then,
\begin{align}
\Sigma_{\bfx} &= \begin{pmatrix} \sigma_1 & 0 \\ 0 & \sigma_2 \end{pmatrix}
\begin{pmatrix} 1 & \rho_{1,2} \\ \rho_{1,2} & 1 \end{pmatrix}
\begin{pmatrix} \sigma_1 & 0 \\ 0 & \sigma_2 \end{pmatrix} \notag \\
&= \begin{pmatrix} \sigma_1 & 0 \\ 0 & \sigma_2 \end{pmatrix} \Omega_{\bfx}
\begin{pmatrix} \sigma_1 & 0 \\ 0 & \sigma_2 \end{pmatrix}; \\
\Omega_{\bfx} &= \begin{pmatrix} \frac{1}{\sigma_1} & 0 \\ 0 & \frac{1}{\sigma_2} \end{pmatrix}
\Sigma_{\bfx} \begin{pmatrix} \frac{1}{\sigma_1} & 0 \\ 0 & \frac{1}{\sigma_2} \end{pmatrix}.
\end{align}
Note that (7.20) and (7.21) are the special cases for $2\times 2$ matrices of (7.16)
and (7.17), respectively.

% NOTE: source literally reads "(7.15) and (7.15)" here (verified on the scan,
% folio 196) -- apparently should reference (7.14) for the covariance matrix.
% Reproduced verbatim per the "never correct the author" rule.
Let $\Sigma_{\bfx}$ and $\Omega_{\bfx}$ be the covariance matrix and the correlation
matrix of three random variables $X_1$, $X_2$, $X_3$ given by (7.15) and (7.15),
respectively. Then,
\begin{align}
\Sigma_{\bfx} &= \begin{pmatrix} \sigma_1 & 0 & 0 \\ 0 & \sigma_2 & 0 \\ 0 & 0 & \sigma_3 \end{pmatrix}
\begin{pmatrix} 1 & \rho_{1,2} & \rho_{1,3} \\ \rho_{1,2} & 1 & \rho_{2,3} \\ \rho_{1,3} & \rho_{2,3} & 1 \end{pmatrix}
\begin{pmatrix} \sigma_1 & 0 & 0 \\ 0 & \sigma_2 & 0 \\ 0 & 0 & \sigma_3 \end{pmatrix} \notag \\
&= \begin{pmatrix} \sigma_1 & 0 & 0 \\ 0 & \sigma_2 & 0 \\ 0 & 0 & \sigma_3 \end{pmatrix}
\Omega_{\bfx}
\begin{pmatrix} \sigma_1 & 0 & 0 \\ 0 & \sigma_2 & 0 \\ 0 & 0 & \sigma_3 \end{pmatrix}; \\
\Omega_{\bfx} &= \begin{pmatrix} \frac{1}{\sigma_1} & 0 & 0 \\ 0 & \frac{1}{\sigma_2} & 0 \\ 0 & 0 & \frac{1}{\sigma_3} \end{pmatrix}
\Sigma_{\bfx}
\begin{pmatrix} \frac{1}{\sigma_1} & 0 & 0 \\ 0 & \frac{1}{\sigma_2} & 0 \\ 0 & 0 & \frac{1}{\sigma_3} \end{pmatrix}.
\end{align}
Note that (7.22) and (7.23) are the special cases for $3\times 3$ matrices of (7.16)
and (7.17), respectively.

% NOTE: The block below prints "Examples:" (plural) as an unnumbered italic run-in
% label in the source, not "Example:" -- so it is NOT the \examplex environment
% (which hard-codes the singular "Example:" text). Closed manually with a plain
% square placed directly after the final display, matching the source, which is
% NOT a \qed/\qedhere marker (no proof is open here).
\par\medskip\noindent\textit{Examples:}\ (i) Consider three random variables $X_1$,
$X_2$, $X_3$ with covariance matrix
\[
\Sigma_{\bfx} \;=\; \begin{pmatrix} 9 & 2 & -6 \\ 2 & 4 & -1 \\ -6 & -1 & 16 \end{pmatrix}.
\]
The standard deviations of $X_1$, $X_2$, and $X_3$ are
\[
\sigma_1 = \sqrt{\Sigma_{\bfx}(1,1)} = 3, \quad \sigma_2 = \sqrt{\Sigma_{\bfx}(2,2)} = 2,
\quad \text{and} \quad \sigma_3 = \sqrt{\Sigma_{\bfx}(3,3)} = 4,
\]
respectively. From (7.23), we find that the correlation matrix of $X_1$, $X_2$, $X_3$
is
\begin{align*}
\Omega_{\bfx} &= \begin{pmatrix} \frac13 & 0 & 0 \\ 0 & \frac12 & 0 \\ 0 & 0 & \frac14 \end{pmatrix}
\begin{pmatrix} 9 & 2 & -6 \\ 2 & 4 & -1 \\ -6 & -1 & 16 \end{pmatrix}
\begin{pmatrix} \frac13 & 0 & 0 \\ 0 & \frac12 & 0 \\ 0 & 0 & \frac14 \end{pmatrix} \\
&= \begin{pmatrix} 1 & \frac13 & -\frac12 \\ \frac13 & 1 & -\frac18 \\ -\frac12 & -\frac18 & 1 \end{pmatrix}.
\end{align*}

(ii) Consider three random variables $X_1$, $X_2$, $X_3$ with correlation matrix
\[
\Omega_{\bfx} \;=\; \begin{pmatrix} 1 & 0.5 & -0.25 \\ 0.5 & 1 & 0.25 \\ -0.25 & 0.25 & 1 \end{pmatrix},
\]
and with standard deviations $\sigma_1 = 4$, $\sigma_2 = 1$, and $\sigma_3 = 6$. From
(7.22), we find that the covariance matrix of $X_1$, $X_2$, $X_3$ is
\begin{align*}
\Sigma_{\bfx} &= \begin{pmatrix} 4 & 0 & 0 \\ 0 & 1 & 0 \\ 0 & 0 & 6 \end{pmatrix}
\begin{pmatrix} 1 & 0.5 & -0.25 \\ 0.5 & 1 & 0.25 \\ -0.25 & 0.25 & 1 \end{pmatrix}
\begin{pmatrix} 4 & 0 & 0 \\ 0 & 1 & 0 \\ 0 & 0 & 6 \end{pmatrix} \\
&= \begin{pmatrix} 16 & 2 & -6 \\ 2 & 1 & 1.5 \\ -6 & 1.5 & 36 \end{pmatrix}. \quad \square
\end{align*}
\par\medskip
